% =============================================================================
% Project : THE_PHILOSOPHY_OF_VIABILITY
% Book    : Book I — Ontology of Viability
% Book ID : VS.B01
%
% File    : 4_ASSETS/tables/
%           VS.B01.P01.C01.S05.SS02.P001.TL01.identity_difference_distinction_matrix.tex
%
% Asset   :
%   Identity, Difference, and Distinction Matrix
%
% Asset ID:
%   VS.B01.P01.C01.S05.SS02.P001.TL01
%
% Purpose :
%   Clarifies the analytical scope of identity, difference,
%   distinction, conceptual differentiation, qualitative
%   difference, and numerical non-identity, while indicating
%   what cannot be inferred from each concept alone.
%
% =============================================================================

\begin{table}[htbp]
\centering
\small
\caption{Identity, Difference, and Distinction: Conceptual and Ontological Scope}
\label{tab:identity-difference-distinction-matrix}

\begin{tabular}{|p{0.19\textwidth}|p{0.35\textwidth}|p{0.35\textwidth}|}
\hline
\textbf{Concept}
&
\textbf{Analytical Scope}
&
\textbf{What It Does Not Establish by Itself}
\\
\hline

Identity
&
The question of whether an existent is itself and what constitutes its being the existent that it is.
&
It is not exhausted by a list of differences, qualities, or distinguishing characteristics.
\\
\hline

Difference
&
Non-sameness between entities, determinations, or conditions in one or more respects.
&
It does not establish that the difference is constitutive of identity or that the entities are numerically distinct.
\\
\hline

Distinction
&
The differentiation of entities, aspects, determinations, or configurations within an ontological or analytical context.
&
It does not necessarily establish independent existence or a plurality of identities.
\\
\hline

Conceptual Differentiation
&
The analytical separation of meanings, descriptions, concepts, or aspects for purposes of examination.
&
It does not establish a corresponding ontological separation or the existence of independent entities.
\\
\hline

Qualitative Difference
&
Variation in characteristics, properties, conditions, or constitution.
&
It does not necessarily establish numerical non-identity or a change of identity.
\\
\hline

Numerical Non-Identity
&
The condition in which two entities are not one and the same existent.
&
It does not provide a positive account of what constitutes the identity of either entity.
\\
\hline

\end{tabular}
\end{table}

% =============================================================================
% END OF TABLE TL01
% =============================================================================